% ===========================================================================
% INTRODUCTION
% File: sections/01_introduction.tex
% ===========================================================================

\section{Introduction}
\label{sec:intro}

% ---------------------------------------------------------------------------
\subsection{The Global Challenge: Food, Energy, and Water Under Climate Pressure}
% ---------------------------------------------------------------------------

\hspace{1cm}By 2050, global food production must increase by approximately
50\% to feed a projected population of nearly ten billion, while freshwater
withdrawals for agriculture --- already accounting for 70\% of the world's
accessible surface water --- are expected to intensify
further~\citep{fao2017, scanlon2023}. In semi-arid and arid regions, which
host some of the fastest-growing populations on Earth, this pressure is
compounded by a pre-existing structural deficit: annual potential
evapotranspiration exceeds precipitation, making every millimetre of saved
irrigation water both economically and ecologically significant.

At the same time, the global energy transition is accelerating the
deployment of ground-mounted solar photovoltaic (PV) installations.
Solar PV has become the least-cost source of new electricity generation
across large parts of the semi-arid belt --- from Southern Europe and the
Middle East to Sub-Saharan Africa and Central Asia~\citep{irena2022}.
However, ground-mounted installations require productive agricultural land,
and the spatial overlap between high-irradiance zones and high-value
irrigated cropland is not coincidental: both are concentrated in the
same semi-arid river basins and coastal plains~\citep{dinesh2016,
ketzer2020}. The result is a land-use conflict that threatens to replace
one resource challenge with another.

% ---------------------------------------------------------------------------
\subsection{Agrivoltaics: Dissolving the Land-Use Conflict}
% ---------------------------------------------------------------------------

\hspace{1cm}Agrivoltaics --- the deliberate co-location of photovoltaic
modules and agricultural crops on the same land parcel --- was proposed
precisely to dissolve this tension. Conceived by \citet{goetzberger1982}
and experimentally validated by \citet{dupraz2011}, the approach rests on
the observation that partial shade from elevated PV panels can simultaneously
generate electricity, reduce crop heat stress, curb soil moisture evaporation,
and increase water-use efficiency for a wide range of
species~\citep{marrou2013, barron2019, elamri2018, semeraro2024}.

The standard metric for quantifying this dual benefit is the Land Equivalent
Ratio (LER), which measures how much land under separate monocultures would
be required to produce the same combined output as an integrated
system~\citep{mead1980}. Published agrivoltaic LER values range from 1.2 to
1.94, with the highest values consistently associated with bifacial modules,
crop-specific panel geometry, and shade-tolerant
species~\citep{riaz2022, mazzeo2025}. A two-component LER of 1.94 means
the integrated system uses land 94\% more efficiently than two separate
operations.

Despite this progress, agrivoltaic research has a systematic blind spot:
the organic residue stream. In a typical tomato production system, crop
residues --- stems, leaves, non-marketable fruit --- amount to approximately
80\% of the harvested fresh mass~\citep{mohammedi2023}. When returned to the
soil, this biomass undergoes aerobic decomposition that can release methane
and nitrous oxide. When combusted, the embodied energy is lost. Neither
outcome is consistent with the circular economy ambitions that motivate
agrivoltaic deployment in the first place.

% ---------------------------------------------------------------------------
\subsection{Anaerobic Digestion: Closing the Residue Loop}
% ---------------------------------------------------------------------------

\hspace{1cm}Anaerobic digestion (AD) converts organic biomass into biogas
and digestate, a nutrient-rich bio-fertilizer. For agricultural systems,
on-farm AD offers three interconnected benefits: it valorizes residues that
would otherwise be wasted, provides dispatchable energy that can offset grid
imports, and returns nitrogen, phosphorus, and potassium to the soil in
plant-available form~\citep{pilarski2025, lallement2023}.

The strategic importance of agricultural biomethane is substantial.
The EU REPowerEU Plan targets 35\,billion m$^3$ of sustainable biomethane
per year by 2030~\citep{ec2022}, and Germany's EEG 2023 and the EU RED III
provide feed-in premiums where biomethane commands approximately 5.5 times
the electricity purchase price~\citep{nikzad2026, eba2025, rabobank2025}.
In Sub-Saharan Africa, biogas from agricultural residues provides an
alternative to imported LPG and diesel in the same farming communities
that stand to benefit most from agrivoltaics~\citep{iea2025}.

A critical thermal challenge limits AD in temperate climates: maintaining
mesophilic conditions (35--40\,$^{\circ}$C) against winter ambient
temperatures below 0\,$^{\circ}$C requires a reliable heat supply.
Conventionally, up to 30\% of biogas is combusted for self-heating,
reducing net energy yield~\citep{lindorfer2008}. PV electricity, abundant
in summer when heating demand is lowest, is a natural complement to
this requirement.

% ---------------------------------------------------------------------------
\subsection{The Missing Integration: APV Meets AD}
% ---------------------------------------------------------------------------

\hspace{1cm}The co-location of agrivoltaics and anaerobic digestion is
logically compelling: APV produces electricity and shade; the crop produces
residues; the digester converts residues to biogas and digestate; digestate
fertilizes the crop. The AD footprint is small and can be sited beneath or
adjacent to the APV array at no additional land cost.

Yet the two technologies have been studied in isolation. Agrivoltaic studies
focus on PV yield and crop PAR, ignoring the residue
stream~\citep{riaz2022, zainali2023, mazzeo2025}. AD studies that include
a PV component treat the solar array as a generic electricity source without
modeling agrivoltaic geometry or crop yield
effects~\citep{pal2024, temiz2024}. The first study to couple APV and AD
optimization was ~\citet{nikzad2026}, limited to one Mediterranean site
without crop yield modeling or a land-use efficiency metric. The consequence
is a systematic gap: no unified metric captures all three land-use outputs,
no optimizer jointly searches both design spaces, and no multi-climate
comparison exists.

% ---------------------------------------------------------------------------
\subsection{Research Gaps}
% ---------------------------------------------------------------------------

\hspace{1cm}Six specific knowledge gaps motivate this study:

\begin{itemize}[leftmargin=*, label=\textbf{G\arabic*:}]

    \item \textbf{No unified land-use metric for integrated APV--AD.}
    The classical two-component LER omits biogas energy. A three-component
    metric is needed to capture the full multi-output land productivity.

    \item \textbf{No joint design optimization.}
    APV geometry and AD design are coupled through crop yield, residue
    production, and digester heating. Separate optimization misses
    cross-variable interactions that determine system-level efficiency.

    \item \textbf{No multi-climate comparative analysis.}
    The sensitivity of the APV--AD synergy to solar resource, temperature,
    and growing-season length has not been systematically characterized
    across contrasting climate zones.

    \item \textbf{No thermal coupling characterization across the
    continental temperature range.}
    Arrhenius BMP suppression ranges from negligible (tropical semi-arid)
    to severe (cold temperate). This variability has not been quantified
    across a four-zone study.

    \item \textbf{No corrected water savings metric.}
    Prior studies report water savings as a fraction of annual ET$_0$,
    systematically underestimating the agronomic benefit by diluting
    growing-season savings with off-season months.

    \item \textbf{No modelling of the AD-to-crop nutrient feedback.}
    The digestate-to-soil loop is absent from all existing agrivoltaic
    studies, preventing a true circular economy assessment.

\end{itemize}

This paper addresses gaps G1--G5. Gap G6 is deferred to future
experimental work requiring multi-season field data from co-located
APV--AD installations.

% ---------------------------------------------------------------------------
\subsection{Contributions of This Work}
% ---------------------------------------------------------------------------

\hspace{1cm}The specific, verifiable contributions of this paper are:

\begin{enumerate}[leftmargin=*, label=\textbf{C\arabic*:}]

    \item \textbf{Novel three-component eLER metric} (G1): $eLER =
    LER_{\text{crop}} + LER_{\text{PV}} + LER_{\text{biogas}}$,
    referenced against three independent monocultures using a
    non-circular PV reference (Definition~A).

    \item \textbf{Integrated six-variable Differential Evolution optimizer}
    (G2): simultaneously searches APV geometry and AD design spaces
    with OLR and inter-row clearance constraints.

    \item \textbf{Four-climate-zone comparative analysis} (G3): Konya
    (BSk), Almer\'{i}a (BSh), Ouagadougou (BSh), and Freiburg (Cfb),
    spanning \SI{12}{\degree}N--\SI{48}{\degree}N.

    \item \textbf{PAR-saturation mechanism} (G3, G4): identification
    and quantification of the counter-intuitive non-monotone relationship
    between solar resource and eLER, driven by tomato light saturation
    at \SI{174}{W\,m^{-2}}.

    \item \textbf{4E performance framework} (G3--G5): Energetic,
    Economic, Energo-economic (LCOE, biogas cost), and Environmental
    (CO$_2$ avoidance, carbon credit) metrics computed jointly for
    all four sites and seven valorization scenarios.

    \item \textbf{Corrected water savings metric} (G5):
    $LER_{\text{water}}$ referenced to growing-season ET$_0$,
    correcting the systematic underestimation of up to 8.1$\times$
    found with annual ET$_0$.

\end{enumerate}

% ---------------------------------------------------------------------------
\subsection{Paper Organization}
% ---------------------------------------------------------------------------

\hspace{1cm}Section~\ref{sec:related} reviews the literature on
agrivoltaic systems, AD integration with solar energy, and LER methodology.
Section~\ref{sec:method} presents the six sub-models, the DE optimizer,
and the 4E framework. Section~\ref{sec:results} reports optimized designs,
eLER decomposition, energy and water performance, and scenario economics
for all four sites. Section~\ref{sec:discussion} interprets the
PAR-saturation mechanism, boundary-convergence behavior, and model
limitations. Section~\ref{sec:conclusion} states the conclusions.
Appendix~\ref{app:de} provides the complete DE pseudocode and Appendix~\ref{app:params}  the complete list of fixed parameters.
